%% Separately authored historical derivative following exact source lines 10353--10674.
%% Necessity: the exact source fixes equipotence, Hartogs, order vocabulary, and
%% four Choice equivalents, but it does not contain the system-variant table,
%% additional calibrated Choice forms, dependency catalog, or binary-tree route.
%% Definition: this file is based on historical slices 1--99, 142--143, and
%% 166--470; source consumers are redirected to exact destinations, internal
%% apparatus IDs are omitted, and the proof/status repairs are declared locally.
%% Standing: this is separately authored reader material, not inherited source
%% or a proof-status promotion of an exact result.
%% Nonredundancy: the removed historical middle duplicates exact definitions;
%% the retained slices add the system table and strictly broader Choice material.
%% Use: parts/p06.tex routes this file once after the exact cumulative-hierarchy
%% chapter, and the annex route supplies only its distinct proof items.

% BEGIN HISTORICAL SLICE 1-99
%% parts/p06/cardinals-choice.tex --- Phase-4 P6 cardinality and choice visit
%% Migrated source: v1.1 lines 10559--10654.  New depth:
%% D-P1-001/002/003/004/006.
%% Apparatus: A-059, A-062, A-063, and A-133.
%% Body = statements and reference apparatus only; proofs are in the annex.

\chapter{Cardinals and Choice}
\label{sec:p06-cardinals-choice}
\chaptermark{Cardinals and Choice}

\noindent\textit{Reader note.}  This separately authored chapter follows the
exact hierarchy/cardinal source.  It adds system comparisons, broader Choice
calibrations, dependency apparatus, and the binary-tree compactness route
without replacing the exact Hartogs, equipotence, or Choice statements.

This chapter separates results available in \(\mathsf{ZF}\) from results
that use a stated choice principle.  It migrates the cardinal-and-choice
entries of the verified source and supplies the additional reference entries
required by the depth register.  The Axiom of Choice itself, the Well-ordering
Theorem, cardinal assignment under Choice, and the Ultrafilter Lemma are first
stated in Part~1; in particular, the Ultrafilter Lemma is referenced through
Definition~\ref{def:ultrafilter-lemma}, not restated as another principle.
Standard references for this chapter are
\cite{Jech2003,Jech1973Choice,HowardRubin1998,Kunen2011}.

\section{Set-theory systems at a glance}

Table~\ref{tab:p06-set-theory-variants} fixes the conventions used when the
handbook compares neighboring set theories.  Here \(\mathsf Z\) means the
modern theory obtained from \(\mathsf{ZF}\) by omitting Replacement while
retaining Foundation; historical uses of ``Zermelo set theory'' vary.
Likewise, this table reserves \(\mathsf{NBG}\) for the presentation with
Global Choice and uses \(\mathsf{GB}\) for its choice-free counterpart.
Global Choice for \(\mathsf{MK}\) is convention-dependent and is therefore
stated separately rather than hidden in the name.  Sources are
\cite{Jech2003,Barwise1975,Mendelson1997,Kelley1975,Quine1937,Specker1953,
Wilshaw2025ConNF}.
\index{Zermelo--Fraenkel set theory}
\index{Zermelo set theory}
\index{Kripke--Platek set theory}
\index{von Neumann--Bernays--Godel set theory@von Neumann--Bernays--G\"odel set theory}
\index{Morse--Kelley set theory}
\index{New Foundations}

\begin{landscape}
\begingroup\scriptsize
\setlength{\tabcolsep}{2.5pt}
\begin{longtable}{@{}p{.48in}p{.95in}p{1.15in}p{1.05in}p{.70in}
  p{.85in}p{1.65in}@{}}
\caption{Set-theory system variants, under the declared
  \(\mathsf Z\), \(\mathsf{NBG}\), and \(\mathsf{MK}\) conventions above.}
\label{tab:p06-set-theory-variants}\\
\toprule
System & Ontology / class sort & Schema strength & Power Set / Replacement &
Foundation & Choice & Comparison \\
\midrule
\endfirsthead
\caption[]{Set-theory system variants, under the declared
  \(\mathsf Z\), \(\mathsf{NBG}\), and \(\mathsf{MK}\) conventions above (continued)}\\
\toprule
System & Ontology / class sort & Schema strength & Power Set / Replacement &
Foundation & Choice & Comparison \\
\midrule
\endhead
\midrule\multicolumn{7}{r}{\emph{continued on next page}}\\
\endfoot
\endlastfoot
\(\mathsf{ZF}\) & One set sort; classes are metalinguistic abbreviations &
Full Separation and Replacement schemata & Power Set: yes; Replacement: yes &
Axiom & Absent; \(\mathsf{ZFC}=\mathsf{ZF}+\mathsf{AC}\) & Baseline
first-order set theory used by the choice calibrations here. \\
\(\mathsf Z\) & One set sort & Full Separation; no Replacement schema &
Power Set: yes; Replacement: no & Retained & Absent;
\(\mathsf{ZC}=\mathsf Z+\mathsf{AC}\) & Declared subtheory of \(\mathsf{ZF}\);
Replacement-dependent constructions are unavailable. \\
\(\mathsf{KP}\) & One set sort & Bounded-formula Separation and
bounded-formula Collection & No Power Set and no full Replacement;
\(\mathsf{KP}^{\mathcal P}\) is separate & Set Induction & Absent in standard
KP & Transitive KP models are admissible; KP is much weaker than
\(\mathsf{ZF}\). \\
\(\mathsf{NBG}\) & Set and class sorts & Predicative elementary class
comprehension & Sets satisfy Power Set and Replacement & For sets & Global
Choice & Conservative over \(\mathsf{ZFC}\) for set-language sentences;
\(\mathsf{GB}\) names the choice-free version. \\
\(\mathsf{MK}\) / \(\mathsf{KM}\) & Set and class sorts & Impredicative class
comprehension permits class quantifiers & Sets satisfy Power Set and
Replacement; stronger class comprehension/collection & For sets & Separate,
convention-dependent addition & Stronger and nonconservative over
\(\mathsf{NBG}\)/\(\mathsf{ZFC}\); no identity with ``ZFC plus one
inaccessible'' is asserted. \\
\(\mathsf{NF}\) & One set sort; no primitive class sort & Stratified
Comprehension & Power sets exist; no unrestricted Replacement & Fails &
Refutable & Neither extension nor subtheory of \(\mathsf{ZF}\); the current
Lean/TTT and paper-level bridge are separated below. \\
\bottomrule
\end{longtable}
\endgroup
\end{landscape}

The 2025 NF project formally verifies in Lean a model construction for
\emph{Tangled Type Theory} (\(\mathsf{TTT}\)); the conclusion
\(\mathsf{Con}(\mathsf{NF})\) then uses the established paper-level
\(\mathsf{NF}\)--\(\mathsf{TTT}\) reduction.  This distinguishes the
machine-checked model construction from the paper-level bridge and from
external peer review
\cite{Wilshaw2025ConNF}.  NF's failure of Choice is the theorem of
\cite{Specker1953}, not a convention of the table.
\index{Tangled Type Theory}

% END HISTORICAL SLICE 1-99

% BEGIN HISTORICAL SLICE 142-143
\section[Partial orders and Choice]{Partial orders and equivalent forms of Choice}

% END HISTORICAL SLICE 142-143

% BEGIN HISTORICAL SLICE 166-470
\begin{hbdefinition}{def:p06-finite-character}{Finite-character family}
Let \(X\) be a set.  A nonempty family
\(\mathcal A\subseteq\mathcal P(X)\) has
\emph{finite character} iff, for every \(Y\subseteq X\),
\[
  Y\in\mathcal A
  \quad\Leftrightarrow\quad
  \text{every finite subset of \(Y\) belongs to \(\mathcal A\)}.
\]
This is a property of a family of subsets, not a finiteness condition on its
members.  It is used in Tukey's formulation of Choice below.
\index{finite character}
\end{hbdefinition}

\begin{hbtheorem}{thm:p06-zorn}{cited=Jech1973Choice}{Choice equivalents and Zorn's lemma}
Over \(\mathsf{ZF}\), the following principles are equivalent.
\begin{enumerate}
  \item The Axiom of Choice as stated in
        Section~\ref{sec:set-prerequisites-06}.
  \item Every set can be well-ordered.
  \item \emph{Zorn's lemma}: if every chain in a poset \(P\), including the
        empty chain, has an upper bound in \(P\), then \(P\) has a maximal
        element.
  \item \emph{Hausdorff's maximal principle}: every chain in a poset is
        contained in a maximal chain.
  \item \emph{Tukey's lemma}: every family of finite character has a maximal
        member under inclusion.
  \item The Cartesian product of every set-indexed family of nonempty sets is
        nonempty.
  \item Tychonoff's theorem for arbitrary compact topological spaces.
\end{enumerate}
The compact-Hausdorff restriction of Tychonoff's theorem is deliberately not
on this list: over \(\mathsf{ZF}\) it is equivalent to the Boolean Prime Ideal
Theorem, hence to the Ultrafilter Lemma, and is strictly weaker than Choice in
the relative-consistency sense stated in
\cref{thm:p01-bpi-choice-separation}.  The equivalence of the Boolean Prime
Ideal Theorem and the Ultrafilter Lemma is proved in
\cref{thm:p01-bpi-equivalence}; its implication from Choice is proved in
\cref{thm:p01-bpi-ac-implication}.
\index{Zorn's lemma}
\index{Hausdorff maximal principle}
\index{Tukey's lemma}
\index{Tychonoff theorem}
\index{Boolean Prime Ideal Theorem}
\end{hbtheorem}

\begin{hbtheorem}{thm:p06-ac-to-zorn}{proved=annex:p6.choice.1}{Choice implies Zorn's lemma}
Assume \(\mathsf{ZF}+\mathsf{AC}\).  If every chain in a poset \(P\),
including the empty chain, has an upper bound in \(P\), then \(P\) has a
maximal element.  The proof uses the well-order of \(P\), transfinite
recursion on the set ordinal \(h(P)\), Replacement for the range at each
stage, and Hartogs's theorem.
\end{hbtheorem}

\begingroup\small
\begin{longtable}{@{}p{.20\textwidth}p{.35\textwidth}p{.36\textwidth}@{}}
\caption{Choice principles and their exact calibration over \(\mathsf{ZF}\).}
\label{tab:p06-ac-equivalents}\\
\toprule
Principle & Form used here & Status over \(\mathsf{ZF}\) \\
\midrule
\endfirsthead
\caption[]{Choice principles and their exact calibration over \(\mathsf{ZF}\) (continued)}\\
\toprule
Principle & Form used here & Status over \(\mathsf{ZF}\) \\
\midrule
\endhead
\midrule\multicolumn{3}{r}{\emph{continued on next page}}\\
\endfoot
\endlastfoot
\multicolumn{3}{@{}l}{\textit{Full-Choice equivalences over \(\mathsf{ZF}\)}}\\*
\addlinespace[2pt]
Axiom of Choice & Choice function for every set of nonempty sets & Reference
form; equivalent to each of the next six rows. \\
Well-ordering theorem & Every set admits a well-order & Equivalent to Choice;
Choice implies the well-ordering statement in
Section~\ref{sec:set-prerequisites-06}, and the converse selects the least
member of each nonempty set. \\
Zorn's lemma & Every chain-upper-bounded poset has a maximal element &
Equivalent to Choice; the complete Choice-to-Zorn direction is
Theorem~\ref{thm:p06-ac-to-zorn} and
Annex~\ref{annex:p6.choice.1}. \\
Hausdorff maximal principle & Every chain extends to a maximal chain &
Equivalent to Choice; applying Zorn to the poset of chains gives the forward
implication. \\
Tukey's lemma & Every finite-character family has an inclusion-maximal member
& Equivalent to Choice; it specializes to maximal partial choices and follows
from Zorn. \\
Product principle & Every set-indexed product of nonempty sets is nonempty &
Equivalent to Choice: a point of the product is exactly a simultaneous choice
function. \\
Tychonoff, arbitrary compact spaces & Every product of compact spaces is
compact & Equivalent to Choice. \\
\addlinespace[4pt]
\multicolumn{3}{@{}l}{\textit{BPI-equivalent and restricted compactness forms}}\\*
\addlinespace[2pt]
Unrestricted propositional compactness & Every finitely satisfiable set of
formulas over an arbitrary propositional variable set is satisfiable &
Equivalent over \(\mathsf{ZF}\) to \(\mathsf{BPI}\); see
Theorem~\ref{thm:p02-arbitrary-prop-compactness-bpi}. \\
Tychonoff, compact Hausdorff spaces & Every product of compact Hausdorff spaces
is compact & Equivalent to the Boolean Prime Ideal Theorem, not to full
Choice. \\
Boolean Prime Ideal Theorem / Ultrafilter Lemma & Every proper Boolean filter
extends to an ultrafilter & Choice implies it; the converse is not provable
over \(\mathsf{ZF}\) if \(\mathsf{ZF}\) is consistent
\cite{HalpernLevy1971,HowardRubin1998}. \\
\bottomrule
\end{longtable}
\endgroup

\begin{hbdefinition}{def:p06-size-comparison}{Comparison of arbitrary set sizes}
For arbitrary sets \(A\) and \(B\), say that the size of \(A\) is
\emph{at most} the size of \(B\) iff an injection \(A\to B\) exists.  Say
that \(A\) has \emph{strictly smaller size} than \(B\) iff an injection
\(A\to B\) exists but no injection \(B\to A\) exists.  These are
metalinguistic relations among sets; they do not presuppose that either set
has an ordinal cardinal.  They differ from equipotence, which requires a
bijection, and the strict relation is used in the trichotomy theorem below.
\index{cardinality!comparison by injection}
\end{hbdefinition}

\begin{hbtheorem}{thm:p06-cardinal-trichotomy}{cited=Jech1973Choice}{Cardinal trichotomy and Choice}
Over \(\mathsf{ZF}\), the following are equivalent: the Axiom of Choice; and,
for every two sets \(A,B\), exactly one of these alternatives holds:
\(A\) has strictly smaller size than \(B\), \(A\approx B\), or \(B\) has
strictly smaller size than \(A\), with comparison understood as in
Definition~\ref{def:p06-size-comparison}.  This is a theorem about
arbitrary sets.  The trichotomy theorem for ordinals requires no Choice.
\end{hbtheorem}

\section[Weaker choice principles]{Weaker choice principles and handbook dependencies}

\begin{hbdefinition}{def:p06-countable-dependent-choice}{Countable Choice and Dependent Choice}
The \emph{Axiom of Countable Choice}, \(\mathsf{AC}_{\omega}\), says that
for every sequence \(\langle X_n:n\in\N\rangle\) of nonempty sets there is
a function \(f:\N\to\bigcup_{n\in\N}X_n\) with \(f(n)\in X_n\) for every
\(n\).  Its restriction to sequences of nonempty finite sets is the
\emph{Axiom of Countable Choice for finite sets},
\(\mathsf{AC}_{\omega}^{\mathrm{fin}}\).  The \emph{Axiom of Dependent
Choice}, \(\mathsf{DC}\), says that whenever \(X\) is nonempty,
\(R\subseteq X\times X\) is serial, and \(x_0\in X\), there is a sequence
\(\langle x_n:n\in\N\rangle\) beginning with \(x_0\) and satisfying
\(x_nRx_{n+1}\) for every \(n\).
These are metatheoretic principles about constructing sequences, not rules of
any object-language calculus.  Dependent Choice is stronger because its next
choice may depend on the preceding value; this distinction is needed in the
dependency table.
\index{Axiom of Countable Choice}
\index{Axiom of Countable Choice!for finite sets}
\index{Dependent Choice}
\end{hbdefinition}

Write \(\mathsf{BPI}\) for the \emph{Boolean Prime Ideal Theorem}, equivalently
the Ultrafilter Lemma of Definition~\ref{def:ultrafilter-lemma}.  This is a
metalinguistic abbreviation for the weaker choice principle already stated in
Part~1; it is used below to keep the implication diagram and dependency table
legible.
\index{Boolean Prime Ideal Theorem}

\begin{hbtheorem}{thm:p06-choice-implications}{cited=HowardRubin1998}{Implications among weaker choice principles}
Over \(\mathsf{ZF}\),
\[
  \begin{aligned}
    \mathsf{AC}&\ \text{ implies }\mathsf{DC},\text{ which implies }
      \mathsf{AC}_{\omega},\text{ which implies }
      \mathsf{AC}_{\omega}^{\mathrm{fin}};\\
    \mathsf{AC}&\ \text{ implies }\mathsf{BPI},\text{ which implies }
      \mathsf{AC}_{\omega}^{\mathrm{fin}}.
  \end{aligned}
\]
The converses along the first chain and the converse from \(\mathsf{BPI}\) to
\(\mathsf{AC}\) fail in suitable models of \(\mathsf{ZF}\).  No comparison
between \(\mathsf{BPI}\) and either \(\mathsf{DC}\) or
\(\mathsf{AC}_{\omega}\) is asserted, and the reverse implication from
\(\mathsf{AC}_{\omega}^{\mathrm{fin}}\) to \(\mathsf{BPI}\) is not asserted.
\end{hbtheorem}
The failure of the \(\mathsf{BPI}\)-to-Choice converse is due to Halpern and
L\'evy; the remaining calibrations and models are catalogued in
\cite{Jech1973Choice,HowardRubin1998,HalpernLevy1971}.

\begingroup\small
\begin{longtable}{@{}p{.27\textwidth}p{.18\textwidth}p{.46\textwidth}@{}}
\caption{Choice-principle dependencies of handbook results.}
\label{tab:p06-choice-dependencies}\\
\toprule
Handbook result or construction & Principle stated & Exact role \\
\midrule
\endfirsthead
\caption[]{Choice-principle dependencies of handbook results (continued)}\\
\toprule
Handbook result or construction & Principle stated & Exact role \\
\midrule
\endhead
\midrule\multicolumn{3}{r}{\emph{continued on next page}}\\
\endfoot
\endlastfoot
Countable propositional compactness,
Theorem~\ref{thm:prop-compactness} & None & The binary-tree route below is a
\(\mathsf{ZF}\) proof; no family of arbitrary finite choice sets is selected. \\
Unrestricted propositional compactness,
Theorem~\ref{thm:p02-arbitrary-prop-compactness-bpi} & \(\mathsf{BPI}\) &
Equivalent over \(\mathsf{ZF}\) to the Boolean Prime Ideal Theorem. \\
Countable-signature first-order completeness,
Corollary~\ref{cor:fol-completeness-countable} & None & The construction
uses one supplied enumeration and least available witnesses. \\
Arbitrary-signature first-order completeness,
Theorem~\ref{thm:fol-completeness} & \(\mathsf{AC}\) & Well-ordering
supports the transfinite Lindenbaum--Henkin construction. \\
Downward L\"owenheim--Skolem,
Theorem~\ref{thm:downward-ls} & \(\mathsf{AC}\) & The theorem states the
ambient Choice hypothesis used for cardinal bounds and witness selection. \\
Nonprincipal-ultrafilter existence,
Lemma~\ref{lem:nonprincipal-ultrafilter} & \(\mathsf{BPI}\) & The
cofinite filter is extended by the Part-1 Ultrafilter Lemma. \\
Zorn's lemma and unrestricted well-ordering,
Theorem~\ref{thm:p06-zorn} & \(\mathsf{AC}\) & Each is equivalent over
\(\mathsf{ZF}\) to full Choice, not merely to a countable fragment. \\
Arbitrary cardinal comparison,
Theorem~\ref{thm:p06-cardinal-trichotomy} & \(\mathsf{AC}\) & Trichotomy for
all sets is equivalent to Choice; ordinal comparison remains a theorem of
\(\mathsf{ZF}\). \\
\L o\'s's theorem, Theorem~\ref{thm:los} & \(\mathsf{AC}\) & The
existential induction step selects witnesses in the satisfying factors; the
ultrafilter itself is already supplied. \\
Part-1 countable-union bound,
Theorem~\ref{thm:infinite-cardinal-finite-coding}, clause 4 &
\(\mathsf{AC}\) is stated;
\(\mathsf{AC}_{\omega}\) supplies the selection step & From bare existence
of injections \(A_n\to\kappa\), a simultaneous sequence of injections is
needed.  The exact restriction consumed there is Countable Choice for such
sets of injections.  Fixed-countable completeness instead begins with one
specified enumeration and does not use this inference. \\
Arbitrary finite successor choices &
\(\mathsf{AC}_{\omega}^{\mathrm{fin}}\) & This principle supplies a selector
for an arbitrary sequence of nonempty finite successor sets.  The canonical
binary recursion below instead uses the fixed order \(0<1\) and requires no
such selector. \\
Generic filter over an externally countable model & Supplied external
enumeration & Once the dense sets and candidate conditions are enumerated,
least compatible choices are made by ordinary recursion in the declared
metatheory; this is not an internal axiom of the model. \\
Dependent Choice & \(\mathsf{DC}\) & No theorem in this table is silently
assigned this principle; any theorem using a dependent recursion not reducible
to ordinary recursion must state \(\mathsf{DC}\) among its hypotheses. \\
\bottomrule
\end{longtable}
\endgroup

\begin{figure}[!htbp]
\centering
\begin{tikzpicture}[
  node distance=18mm and 15mm,
  principle/.style={draw,rounded corners,minimum width=17mm,minimum height=8mm,
                    inner sep=3pt,font=\small},
  implication/.style={-{Latex[length=2.2mm]},thick},
  note/.style={font=\scriptsize,align=center}
]
  \node[principle] (ac) {\(\mathsf{AC}\)};
  \node[principle,right=of ac] (dc) {\(\mathsf{DC}\)};
  \node[principle,right=of dc] (acw) {\(\mathsf{AC}_{\omega}\)};
  \node[principle,right=of acw] (acwf) {\(\mathsf{AC}_{\omega}^{\mathrm{fin}}\)};
  \node[principle,below=of dc] (bpi) {\(\mathsf{BPI}\)};
  \draw[implication] (ac) -- (dc);
  \draw[implication] (dc) -- (acw);
  \draw[implication] (acw) -- (acwf);
  \draw[implication] (ac) |- (bpi);
  \draw[implication] (bpi) -| (acwf);
  \node[note,above=1mm of dc] {converse fails};
  \node[note,above=1mm of acw] {converse fails};
  \node[note,above=1mm of acwf] {converse fails};
  \node[note,left=2mm of bpi] {converse\\fails};
  \node[note,below=3mm of bpi,anchor=north west] {No comparison with
    \(\mathsf{DC}\) or \(\mathsf{AC}_{\omega}\);\\
    the reverse diagonal is not asserted.};
\end{tikzpicture}
\caption{Implications among the choice principles used in the handbook.  Every
arrow is provable in \(\mathsf{ZF}\); the marked
  converses fail in suitable models, while the reverse diagonal is left
  unasserted.}
\label{fig:p06-choice-implications}
\end{figure}
\clearpage

\section[Binary-tree compactness]{Binary trees, the fan theorem, and compactness}

\begin{hbdefinition}{def:p06-binary-tree-bar}{Binary trees, branches, and bars}
Write \(2^{<\omega}\) for the finite binary sequences and \(2^\omega\) for
the infinite binary sequences.  A \emph{binary tree} is a prefix-closed set
\(\mathcal T\subseteq2^{<\omega}\).  It has nodes at every level iff for every
\(n\in\N\) some \(\mathbf s\in\mathcal T\) has length \(n\).  A \emph{branch} through
\(\mathcal T\) is an \(f\in2^\omega\) with every finite initial segment
\(f\mathord{\upharpoonright}n\in\mathcal T\).  A \emph{bar} is a set
\(B\subseteq2^{<\omega}\) such that every \(f\in2^\omega\) has an initial
segment in \(B\).  These notions concern the canonical binary tree; replacing
it by an arbitrary sequence of nonempty finite successor sets can introduce a
use of \(\mathsf{AC}_{\omega}^{\mathrm{fin}}\).
\index{binary tree}
\index{branch!through a binary tree}
\index{bar!on the binary tree}
\end{hbdefinition}

\begin{hbtheorem}{thm:p06-konig-infinity}{proved=annex:p6.choice.2}{K\"onig's infinity lemma for the binary tree}
In \(\mathsf{ZF}\), if a binary tree \(\mathcal T\subseteq2^{<\omega}\) has a
node at every level, then \(\mathcal T\) has an infinite branch.  No choice
principle is used.
\end{hbtheorem}
The canonical binary formulation is essential to the choice-free claim; see
\cite{Jech2003} for the surrounding set-theoretic compactness methods.

\begin{hbcorollary}{thm:p06-classical-fan}{proved=annex:p6.choice.3}{Classical fan theorem}
In \(\mathsf{ZF}\), if \(B\subseteq2^{<\omega}\) is a bar, then there is an
\(N\in\N\) such that every \(\mathbf s\in2^N\) has an initial segment in \(B\).
This is the classical contrapositive form supplied by K\"onig's infinity
lemma, not a claim about constructive fan principles.
\end{hbcorollary}

\begin{hbmethod}{def:p06-countable-compactness-route}{Choice-free route to countable propositional compactness}
\begin{description}
  \item[Input.] A finitely satisfiable set \(\Gamma\) of formulas over a fixed
        countable propositional alphabet, together with enumerations of its
        variables and formulas.
  \item[Tree.] At level \(n\), retain exactly the binary assignments to the
        first \(n\) variables that extend to a valuation satisfying the first
        \(n\) formulas.
  \item[Dependency.] Apply Theorem~\ref{thm:p06-konig-infinity} to this
        canonical binary tree; read the resulting branch as one valuation of
        the whole alphabet.
  \item[Result.] The correctness statement is the already proved
        Theorem~\ref{thm:prop-compactness}.  For an unrestricted variable set,
        Theorem~\ref{thm:p02-arbitrary-prop-compactness-bpi} records the exact
        \(\mathsf{BPI}\) calibration.
\end{description}
\index{propositional compactness!choice-free route}
\end{hbmethod}
% END HISTORICAL SLICE 166-470
